
\documentclass{article}
\usepackage{colortbl}
\usepackage{makecell}
\usepackage{multirow}
\usepackage{supertabular}

\begin{document}

\newcounter{utterance}

\centering \large Interaction Transcript for game `codenames', experiment `association\_easy', episode 3 with glm{-}5.2.
\vspace{24pt}

{ \footnotesize  \setcounter{utterance}{1}
\setlength{\tabcolsep}{0pt}
\begin{supertabular}{c@{$\;$}|p{.15\linewidth}@{}p{.15\linewidth}p{.15\linewidth}p{.15\linewidth}p{.15\linewidth}p{.15\linewidth}}
    \# & $\;$A & \multicolumn{4}{c}{Game Master} & $\;\:$B\\
    \hline

    \theutterance \stepcounter{utterance}  
    & & & \multicolumn{2}{p{0.3\linewidth}}{
        \cellcolor[rgb]{0.9,0.9,0.9}{
            \makecell[{{p{\linewidth}}}]{
                \texttt{\tiny{[GM$|$GM]}}
                {'hidden': {'team': ['vītītis', 'viesuļvētra', 'krāpšana', 'laupīšana', 'ātruma pārsniegšana', 'lietus', 'kardināls', 'tornado', 'taifūns'], 'innocent': ['pēda', 'krekls', 'jards', 'kaklasaite', 'sporta krekls', 'kilometrs', 'kostīms'], 'opponent': ['kliņģerīte', 'roze', 'grāmatvedis', 'pistole', 'tops', 'ģerānija', 'jurists', 'mehāniķis'], 'assassin': ['lielgabals']}, 'revealed': {'team': {'team': [], 'innocent': [], 'opponent': [], 'assassin': []}, 'opponent': {'team': [], 'innocent': [], 'opponent': [], 'assassin': []}}}
            }
        }
    }
    & & \\ \\

    \theutterance \stepcounter{utterance}  
    & & \multicolumn{4}{p{0.6\linewidth}}{
        \cellcolor[rgb]{0.9,0.9,0.9}{
            \makecell[{{p{\linewidth}}}]{
                \texttt{\tiny{[P1$\langle$GM]}}
                \texttt{Sniedziet vienu vienvārdu norādi, kas ir saistīta ar kādu no jūsu komandas vārdiem no tālāk dotā saraksta. Jūs varat izvēlēties, uz cik un kādiem vārdiem attiecas jūsu norāde.} \\
\texttt{Mērķis ir panākt, lai jūsu partneris uzmin visus jūsu komandas vārdus, pamatojoties uz jūsu dotajām norādēm, pirms otra komanda ir atklājusi visus savus vārdus. Ja jūsu partneris uzmin "pretinieka" vai "novirzītāja" vārdu, viņa uzminēšanas kārta beidzas, un pretinieku komanda var veikt savu kārtu. Ja jūsu partneris uzmin "nogalināšanas" vārdu, spēle beidzas un jūsu komanda uzreiz zaudē.} \\
\texttt{Jo vairāk vārdu, uz kuriem attiecas jūsu norāde, jo ātrāk jūs varat uzvarēt spēli, bet jo grūtāk asociācijas var būt uzminēt jūsu partnerim, tāpēc izvēlieties savu norādes vārdu gudri. Norādes vārdam ir jābūt semantiski saistītam ar mērķa vārdiem, tas nedrīkst būt viens no sarakstos iekļautajiem vārdiem vai saturēt to daļas.} \\
\\ 
\texttt{Vienvārdu norādi un ar komatu atdalīto saistīto mērķa vārdu sarakstu vienmēr sniedziet šādā formātā un atbildi veidojiet pēc iespējas īsāku, nekad neiekļaujiet citu tekstu, kas nav prasīts šajā formā:} \\
\\ 
\texttt{PADOMS: <VĀRDS>} \\
\texttt{MĒRĶI: <MĒRĶI>} \\
\\ 
\texttt{Jūsu komandas vārdi ir: vītītis, viesuļvētra, krāpšana, laupīšana, ātruma pārsniegšana, lietus, kardināls, tornado, taifūns.} \\
\texttt{Jūsu pretinieka komandas vārdi ir: kliņģerīte, roze, grāmatvedis, pistole, tops, ģerānija, jurists, mehāniķis.} \\
\texttt{Novirzītāja vārdi ir: pēda, krekls, jards, kaklasaite, sporta krekls, kilometrs, kostīms.} \\
\texttt{Nogalināšanas vārdi ir: lielgabals.} \\
            }
        }
    }
    & \\ \\

    \theutterance \stepcounter{utterance}  
    & \multicolumn{4}{p{0.6\linewidth}}{
        \cellcolor[rgb]{0.9,0.9,0.9}{
            \makecell[{{p{\linewidth}}}]{
                \texttt{\tiny{[P1$\rangle$GM]}}
                \texttt{PADOMS: vējš} \\
\texttt{MĒRĶI: viesuļvētra, tornado, taifūns} \\
            }
        }
    }
    & & \\ \\

    \theutterance \stepcounter{utterance}  
    & & & \multicolumn{2}{p{0.3\linewidth}}{
        \cellcolor[rgb]{0.9,0.9,0.9}{
            \makecell[{{p{\linewidth}}}]{
                \texttt{\tiny{[GM$|$GM]}}
                \texttt{vējš} \\
            }
        }
    }
    & & \\ \\

    \theutterance \stepcounter{utterance}  
    & & & \multicolumn{2}{p{0.3\linewidth}}{
        \cellcolor[rgb]{0.9,0.9,0.9}{
            \makecell[{{p{\linewidth}}}]{
                \texttt{\tiny{[GM$|$GM]}}
                ['viesuļvētra', 'tornado', 'taifūns']
            }
        }
    }
    & & \\ \\

    \theutterance \stepcounter{utterance}  
    & & & \multicolumn{2}{p{0.3\linewidth}}{
        \cellcolor[rgb]{0.9,0.9,0.9}{
            \makecell[{{p{\linewidth}}}]{
                \texttt{\tiny{[GM$|$GM]}}
                {'word': 'viesuļvētra', 'assignment': 'team'}
            }
        }
    }
    & & \\ \\

    \theutterance \stepcounter{utterance}  
    & & & \multicolumn{2}{p{0.3\linewidth}}{
        \cellcolor[rgb]{0.9,0.9,0.9}{
            \makecell[{{p{\linewidth}}}]{
                \texttt{\tiny{[GM$|$GM]}}
                {'word': 'tornado', 'assignment': 'team'}
            }
        }
    }
    & & \\ \\

    \theutterance \stepcounter{utterance}  
    & & & \multicolumn{2}{p{0.3\linewidth}}{
        \cellcolor[rgb]{0.9,0.9,0.9}{
            \makecell[{{p{\linewidth}}}]{
                \texttt{\tiny{[GM$|$GM]}}
                {'word': 'taifūns', 'assignment': 'team'}
            }
        }
    }
    & & \\ \\

    \theutterance \stepcounter{utterance}  
    & & \multicolumn{4}{p{0.6\linewidth}}{
        \cellcolor[rgb]{0.9,0.9,0.9}{
            \makecell[{{p{\linewidth}}}]{
                \texttt{\tiny{[GM$\rangle$P2]}}
                \texttt{Sniedziet ar komatiem atdalītu sarakstu ar līdz 3 vārdiem no turpmāk minētā saraksta, kas vislabāk attiecas uz vārdu 'vējš' vai ir visciešāk saistīti ar to. Vienmēr sāciet minējumu sarakstu ar MINĒJUMS:  un atbildē neiekļaujiet nekādu citu tekstu.} \\
\\ 
\texttt{jurists, laupīšana, lietus, vītītis, tops, roze, mehāniķis, kardināls, kliņģerīte, pistole, taifūns, ģerānija, sporta krekls, kilometrs, kostīms, tornado, lielgabals, kaklasaite, grāmatvedis, krāpšana, jards, ātruma pārsniegšana, krekls, viesuļvētra, pēda} \\
            }
        }
    }
    & \\ \\

    \theutterance \stepcounter{utterance}  
    & & & \multicolumn{4}{p{0.6\linewidth}}{
        \cellcolor[rgb]{0.9,0.9,0.9}{
            \makecell[{{p{\linewidth}}}]{
                \texttt{\tiny{[GM$\langle$P2]}}
                \texttt{MINĒJUMS: taifūns, tornado, viesuļvētra} \\
            }
        }
    }
     \\ \\

    \theutterance \stepcounter{utterance}  
    & & & \multicolumn{2}{p{0.3\linewidth}}{
        \cellcolor[rgb]{0.9,0.9,0.9}{
            \makecell[{{p{\linewidth}}}]{
                \texttt{\tiny{[GM$|$GM]}}
                ['taifūns', 'tornado', 'viesuļvētra']
            }
        }
    }
    & & \\ \\

    \theutterance \stepcounter{utterance}  
    & & & \multicolumn{2}{p{0.3\linewidth}}{
        \cellcolor[rgb]{0.9,0.9,0.9}{
            \makecell[{{p{\linewidth}}}]{
                \texttt{\tiny{[GM$|$GM]}}
                {'word': 'taifūns', 'assignment': 'team'}
            }
        }
    }
    & & \\ \\

    \theutterance \stepcounter{utterance}  
    & & & \multicolumn{2}{p{0.3\linewidth}}{
        \cellcolor[rgb]{0.9,0.9,0.9}{
            \makecell[{{p{\linewidth}}}]{
                \texttt{\tiny{[GM$|$GM]}}
                {'word': 'taifūns', 'assignment': 'team'}
            }
        }
    }
    & & \\ \\

    \theutterance \stepcounter{utterance}  
    & & & \multicolumn{2}{p{0.3\linewidth}}{
        \cellcolor[rgb]{0.9,0.9,0.9}{
            \makecell[{{p{\linewidth}}}]{
                \texttt{\tiny{[GM$|$GM]}}
                {'word': 'tornado', 'assignment': 'team'}
            }
        }
    }
    & & \\ \\

    \theutterance \stepcounter{utterance}  
    & & & \multicolumn{2}{p{0.3\linewidth}}{
        \cellcolor[rgb]{0.9,0.9,0.9}{
            \makecell[{{p{\linewidth}}}]{
                \texttt{\tiny{[GM$|$GM]}}
                {'word': 'tornado', 'assignment': 'team'}
            }
        }
    }
    & & \\ \\

    \theutterance \stepcounter{utterance}  
    & & & \multicolumn{2}{p{0.3\linewidth}}{
        \cellcolor[rgb]{0.9,0.9,0.9}{
            \makecell[{{p{\linewidth}}}]{
                \texttt{\tiny{[GM$|$GM]}}
                {'word': 'viesuļvētra', 'assignment': 'team'}
            }
        }
    }
    & & \\ \\

    \theutterance \stepcounter{utterance}  
    & & & \multicolumn{2}{p{0.3\linewidth}}{
        \cellcolor[rgb]{0.9,0.9,0.9}{
            \makecell[{{p{\linewidth}}}]{
                \texttt{\tiny{[GM$|$GM]}}
                {'word': 'viesuļvētra', 'assignment': 'team'}
            }
        }
    }
    & & \\ \\

    \theutterance \stepcounter{utterance}  
    & & & \multicolumn{2}{p{0.3\linewidth}}{
        \cellcolor[rgb]{0.9,0.9,0.9}{
            \makecell[{{p{\linewidth}}}]{
                \texttt{\tiny{[GM$|$GM]}}
                {'word': 'jurists', 'assignment': 'opponent'}
            }
        }
    }
    & & \\ \\

    \theutterance \stepcounter{utterance}  
    & & & \multicolumn{2}{p{0.3\linewidth}}{
        \cellcolor[rgb]{0.9,0.9,0.9}{
            \makecell[{{p{\linewidth}}}]{
                \texttt{\tiny{[GM$|$GM]}}
                {'hidden': {'team': ['vītītis', 'krāpšana', 'laupīšana', 'ātruma pārsniegšana', 'lietus', 'kardināls'], 'innocent': ['pēda', 'krekls', 'jards', 'kaklasaite', 'sporta krekls', 'kilometrs', 'kostīms'], 'opponent': ['kliņģerīte', 'roze', 'grāmatvedis', 'pistole', 'tops', 'ģerānija', 'mehāniķis'], 'assassin': ['lielgabals']}, 'revealed': {'team': {'team': ['taifūns', 'tornado', 'viesuļvētra'], 'innocent': [], 'opponent': [], 'assassin': []}, 'opponent': {'team': [], 'innocent': [], 'opponent': ['jurists'], 'assassin': []}}}
            }
        }
    }
    & & \\ \\

    \theutterance \stepcounter{utterance}  
    & & \multicolumn{4}{p{0.6\linewidth}}{
        \cellcolor[rgb]{0.9,0.9,0.9}{
            \makecell[{{p{\linewidth}}}]{
                \texttt{\tiny{[P1$\langle$GM]}}
                \texttt{Vārdi taifūns, tornado, viesuļvētra tika uzminēti pareizi. Jūsu komandas biedra gājiens beidzās tur.} \\
\texttt{Tagad sniedziet vēl vienu norādi, kas attiecas uz dažiem no jūsu atlikušajiem komandas vārdiem, un saistīto mērķa vārdu sarakstu. Neaizmirstiet sākt savu norādi ar 'PADOMS: ', pēc tam ievietot jaunu rindu un sākt ar komatiem atdalīto mērķa vārdu sarakstu ar 'MĒRĶI: '. Piezīme: salīdzinājumā ar iepriekšējiem pieprasījumiem daži vārdi no sarakstiem ir izņemti.} \\
\\ 
\texttt{Jūsu atlikušie komandas vārdi ir: vītītis, krāpšana, laupīšana, ātruma pārsniegšana, lietus, kardināls.} \\
\texttt{Atlikušie pretinieka vārdi ir: kliņģerīte, roze, grāmatvedis, pistole, tops, ģerānija, jurists, mehāniķis.} \\
\texttt{Atlikušie traucētājvārdi ir: pēda, krekls, jards, kaklasaite, sporta krekls, kilometrs, kostīms.} \\
\texttt{Atlikušie slepkavības vārdi ir: lielgabals.} \\
            }
        }
    }
    & \\ \\

    \theutterance \stepcounter{utterance}  
    & \multicolumn{4}{p{0.6\linewidth}}{
        \cellcolor[rgb]{0.9,0.9,0.9}{
            \makecell[{{p{\linewidth}}}]{
                \texttt{\tiny{[P1$\rangle$GM]}}
                \texttt{PADOMS: noziegums} \\
\texttt{MĒRĶI: krāpšana, laupīšana} \\
            }
        }
    }
    & & \\ \\

    \theutterance \stepcounter{utterance}  
    & & & \multicolumn{2}{p{0.3\linewidth}}{
        \cellcolor[rgb]{0.9,0.9,0.9}{
            \makecell[{{p{\linewidth}}}]{
                \texttt{\tiny{[GM$|$GM]}}
                \texttt{noziegums} \\
            }
        }
    }
    & & \\ \\

    \theutterance \stepcounter{utterance}  
    & & & \multicolumn{2}{p{0.3\linewidth}}{
        \cellcolor[rgb]{0.9,0.9,0.9}{
            \makecell[{{p{\linewidth}}}]{
                \texttt{\tiny{[GM$|$GM]}}
                ['krāpšana', 'laupīšana']
            }
        }
    }
    & & \\ \\

    \theutterance \stepcounter{utterance}  
    & & & \multicolumn{2}{p{0.3\linewidth}}{
        \cellcolor[rgb]{0.9,0.9,0.9}{
            \makecell[{{p{\linewidth}}}]{
                \texttt{\tiny{[GM$|$GM]}}
                {'word': 'krāpšana', 'assignment': 'team'}
            }
        }
    }
    & & \\ \\

    \theutterance \stepcounter{utterance}  
    & & & \multicolumn{2}{p{0.3\linewidth}}{
        \cellcolor[rgb]{0.9,0.9,0.9}{
            \makecell[{{p{\linewidth}}}]{
                \texttt{\tiny{[GM$|$GM]}}
                {'word': 'laupīšana', 'assignment': 'team'}
            }
        }
    }
    & & \\ \\

    \theutterance \stepcounter{utterance}  
    & & \multicolumn{4}{p{0.6\linewidth}}{
        \cellcolor[rgb]{0.9,0.9,0.9}{
            \makecell[{{p{\linewidth}}}]{
                \texttt{\tiny{[GM$\rangle$P2]}}
                \texttt{Vārdi taifūns, tornado, viesuļvētra tika uzminēti pareizi. Jūsu gājiens beidzās tur.} \\
\texttt{Tagad sniedziet vēl vienu ar komatiem atdalītu sarakstu, kurā ir vismaz 1 un ne vairāk kā 2 vārdi no turpmāk minētā vārdu saraksta, kas vislabāk sasaucas vai ir visciešāk saistīti ar vārdu 'noziegums'. Neaizmirstiet atbildi sākt ar MINĒJUMS: . Piezīme: salīdzinājumā ar iepriekšējiem pieprasījumiem daži vārdi no saraksta ir izņemti.} \\
\\ 
\texttt{laupīšana, lietus, vītītis, tops, roze, mehāniķis, kardināls, kliņģerīte, pistole, ģerānija, sporta krekls, kilometrs, kostīms, lielgabals, kaklasaite, grāmatvedis, krāpšana, jards, ātruma pārsniegšana, krekls, pēda} \\
            }
        }
    }
    & \\ \\

    \theutterance \stepcounter{utterance}  
    & & & \multicolumn{4}{p{0.6\linewidth}}{
        \cellcolor[rgb]{0.9,0.9,0.9}{
            \makecell[{{p{\linewidth}}}]{
                \texttt{\tiny{[GM$\langle$P2]}}
                \texttt{MINĒJUMS: laupīšana, krāpšana} \\
            }
        }
    }
     \\ \\

    \theutterance \stepcounter{utterance}  
    & & & \multicolumn{2}{p{0.3\linewidth}}{
        \cellcolor[rgb]{0.9,0.9,0.9}{
            \makecell[{{p{\linewidth}}}]{
                \texttt{\tiny{[GM$|$GM]}}
                ['laupīšana', 'krāpšana']
            }
        }
    }
    & & \\ \\

    \theutterance \stepcounter{utterance}  
    & & & \multicolumn{2}{p{0.3\linewidth}}{
        \cellcolor[rgb]{0.9,0.9,0.9}{
            \makecell[{{p{\linewidth}}}]{
                \texttt{\tiny{[GM$|$GM]}}
                {'word': 'laupīšana', 'assignment': 'team'}
            }
        }
    }
    & & \\ \\

    \theutterance \stepcounter{utterance}  
    & & & \multicolumn{2}{p{0.3\linewidth}}{
        \cellcolor[rgb]{0.9,0.9,0.9}{
            \makecell[{{p{\linewidth}}}]{
                \texttt{\tiny{[GM$|$GM]}}
                {'word': 'laupīšana', 'assignment': 'team'}
            }
        }
    }
    & & \\ \\

    \theutterance \stepcounter{utterance}  
    & & & \multicolumn{2}{p{0.3\linewidth}}{
        \cellcolor[rgb]{0.9,0.9,0.9}{
            \makecell[{{p{\linewidth}}}]{
                \texttt{\tiny{[GM$|$GM]}}
                {'word': 'krāpšana', 'assignment': 'team'}
            }
        }
    }
    & & \\ \\

    \theutterance \stepcounter{utterance}  
    & & & \multicolumn{2}{p{0.3\linewidth}}{
        \cellcolor[rgb]{0.9,0.9,0.9}{
            \makecell[{{p{\linewidth}}}]{
                \texttt{\tiny{[GM$|$GM]}}
                {'word': 'krāpšana', 'assignment': 'team'}
            }
        }
    }
    & & \\ \\

    \theutterance \stepcounter{utterance}  
    & & & \multicolumn{2}{p{0.3\linewidth}}{
        \cellcolor[rgb]{0.9,0.9,0.9}{
            \makecell[{{p{\linewidth}}}]{
                \texttt{\tiny{[GM$|$GM]}}
                {'word': 'roze', 'assignment': 'opponent'}
            }
        }
    }
    & & \\ \\

    \theutterance \stepcounter{utterance}  
    & & & \multicolumn{2}{p{0.3\linewidth}}{
        \cellcolor[rgb]{0.9,0.9,0.9}{
            \makecell[{{p{\linewidth}}}]{
                \texttt{\tiny{[GM$|$GM]}}
                {'hidden': {'team': ['vītītis', 'ātruma pārsniegšana', 'lietus', 'kardināls'], 'innocent': ['pēda', 'krekls', 'jards', 'kaklasaite', 'sporta krekls', 'kilometrs', 'kostīms'], 'opponent': ['kliņģerīte', 'grāmatvedis', 'pistole', 'tops', 'ģerānija', 'mehāniķis'], 'assassin': ['lielgabals']}, 'revealed': {'team': {'team': ['taifūns', 'tornado', 'viesuļvētra', 'laupīšana', 'krāpšana'], 'innocent': [], 'opponent': [], 'assassin': []}, 'opponent': {'team': [], 'innocent': [], 'opponent': ['jurists', 'roze'], 'assassin': []}}}
            }
        }
    }
    & & \\ \\

    \theutterance \stepcounter{utterance}  
    & & \multicolumn{4}{p{0.6\linewidth}}{
        \cellcolor[rgb]{0.9,0.9,0.9}{
            \makecell[{{p{\linewidth}}}]{
                \texttt{\tiny{[P1$\langle$GM]}}
                \texttt{Vārdi laupīšana, krāpšana tika uzminēti pareizi. Jūsu komandas biedra gājiens beidzās tur.} \\
\texttt{Tagad sniedziet vēl vienu norādi, kas attiecas uz dažiem no jūsu atlikušajiem komandas vārdiem, un saistīto mērķa vārdu sarakstu. Neaizmirstiet sākt savu norādi ar 'PADOMS: ', pēc tam ievietot jaunu rindu un sākt ar komatiem atdalīto mērķa vārdu sarakstu ar 'MĒRĶI: '. Piezīme: salīdzinājumā ar iepriekšējiem pieprasījumiem daži vārdi no sarakstiem ir izņemti.} \\
\\ 
\texttt{Jūsu atlikušie komandas vārdi ir: vītītis, ātruma pārsniegšana, lietus, kardināls.} \\
\texttt{Atlikušie pretinieka vārdi ir: kliņģerīte, roze, grāmatvedis, pistole, tops, ģerānija, mehāniķis.} \\
\texttt{Atlikušie traucētājvārdi ir: pēda, krekls, jards, kaklasaite, sporta krekls, kilometrs, kostīms.} \\
\texttt{Atlikušie slepkavības vārdi ir: lielgabals.} \\
            }
        }
    }
    & \\ \\

    \theutterance \stepcounter{utterance}  
    & \multicolumn{4}{p{0.6\linewidth}}{
        \cellcolor[rgb]{0.9,0.9,0.9}{
            \makecell[{{p{\linewidth}}}]{
                \texttt{\tiny{[P1$\rangle$GM]}}
                \texttt{PADOMS: mērs} \\
\texttt{MĒRĶI: ātruma pārsniegšana} \\
            }
        }
    }
    & & \\ \\

    \theutterance \stepcounter{utterance}  
    & & & \multicolumn{2}{p{0.3\linewidth}}{
        \cellcolor[rgb]{0.9,0.9,0.9}{
            \makecell[{{p{\linewidth}}}]{
                \texttt{\tiny{[GM$|$GM]}}
                \texttt{mērs} \\
            }
        }
    }
    & & \\ \\

    \theutterance \stepcounter{utterance}  
    & & & \multicolumn{2}{p{0.3\linewidth}}{
        \cellcolor[rgb]{0.9,0.9,0.9}{
            \makecell[{{p{\linewidth}}}]{
                \texttt{\tiny{[GM$|$GM]}}
                ['ātruma pārsniegšana']
            }
        }
    }
    & & \\ \\

    \theutterance \stepcounter{utterance}  
    & & & \multicolumn{2}{p{0.3\linewidth}}{
        \cellcolor[rgb]{0.9,0.9,0.9}{
            \makecell[{{p{\linewidth}}}]{
                \texttt{\tiny{[GM$|$GM]}}
                {'word': 'ātruma pārsniegšana', 'assignment': 'team'}
            }
        }
    }
    & & \\ \\

    \theutterance \stepcounter{utterance}  
    & & \multicolumn{4}{p{0.6\linewidth}}{
        \cellcolor[rgb]{0.9,0.9,0.9}{
            \makecell[{{p{\linewidth}}}]{
                \texttt{\tiny{[GM$\rangle$P2]}}
                \texttt{Vārdi laupīšana, krāpšana tika uzminēti pareizi. Jūsu gājiens beidzās tur.} \\
\texttt{Tagad sniedziet vēl vienu ar komatiem atdalītu sarakstu, kurā ir vismaz 1 un ne vairāk kā 1 vārdi no turpmāk minētā vārdu saraksta, kas vislabāk sasaucas vai ir visciešāk saistīti ar vārdu 'mērs'. Neaizmirstiet atbildi sākt ar MINĒJUMS: . Piezīme: salīdzinājumā ar iepriekšējiem pieprasījumiem daži vārdi no saraksta ir izņemti.} \\
\\ 
\texttt{lietus, vītītis, tops, mehāniķis, kardināls, kliņģerīte, pistole, ģerānija, sporta krekls, kilometrs, kostīms, lielgabals, kaklasaite, grāmatvedis, jards, ātruma pārsniegšana, krekls, pēda} \\
            }
        }
    }
    & \\ \\

    \theutterance \stepcounter{utterance}  
    & & & \multicolumn{4}{p{0.6\linewidth}}{
        \cellcolor[rgb]{0.9,0.9,0.9}{
            \makecell[{{p{\linewidth}}}]{
                \texttt{\tiny{[GM$\langle$P2]}}
                \texttt{MINĒJUMS: jards} \\
            }
        }
    }
     \\ \\

    \theutterance \stepcounter{utterance}  
    & & & \multicolumn{2}{p{0.3\linewidth}}{
        \cellcolor[rgb]{0.9,0.9,0.9}{
            \makecell[{{p{\linewidth}}}]{
                \texttt{\tiny{[GM$|$GM]}}
                ['jards']
            }
        }
    }
    & & \\ \\

    \theutterance \stepcounter{utterance}  
    & & & \multicolumn{2}{p{0.3\linewidth}}{
        \cellcolor[rgb]{0.9,0.9,0.9}{
            \makecell[{{p{\linewidth}}}]{
                \texttt{\tiny{[GM$|$GM]}}
                {'word': 'jards', 'assignment': 'innocent'}
            }
        }
    }
    & & \\ \\

    \theutterance \stepcounter{utterance}  
    & & & \multicolumn{2}{p{0.3\linewidth}}{
        \cellcolor[rgb]{0.9,0.9,0.9}{
            \makecell[{{p{\linewidth}}}]{
                \texttt{\tiny{[GM$|$GM]}}
                \texttt{jards} \\
            }
        }
    }
    & & \\ \\

    \theutterance \stepcounter{utterance}  
    & & & \multicolumn{2}{p{0.3\linewidth}}{
        \cellcolor[rgb]{0.9,0.9,0.9}{
            \makecell[{{p{\linewidth}}}]{
                \texttt{\tiny{[GM$|$GM]}}
                {'word': 'tops', 'assignment': 'opponent'}
            }
        }
    }
    & & \\ \\

    \theutterance \stepcounter{utterance}  
    & & & \multicolumn{2}{p{0.3\linewidth}}{
        \cellcolor[rgb]{0.9,0.9,0.9}{
            \makecell[{{p{\linewidth}}}]{
                \texttt{\tiny{[GM$|$GM]}}
                {'hidden': {'team': ['vītītis', 'ātruma pārsniegšana', 'lietus', 'kardināls'], 'innocent': ['pēda', 'krekls', 'kaklasaite', 'sporta krekls', 'kilometrs', 'kostīms'], 'opponent': ['kliņģerīte', 'grāmatvedis', 'pistole', 'ģerānija', 'mehāniķis'], 'assassin': ['lielgabals']}, 'revealed': {'team': {'team': ['taifūns', 'tornado', 'viesuļvētra', 'laupīšana', 'krāpšana'], 'innocent': ['jards'], 'opponent': [], 'assassin': []}, 'opponent': {'team': [], 'innocent': [], 'opponent': ['jurists', 'roze', 'tops'], 'assassin': []}}}
            }
        }
    }
    & & \\ \\

    \theutterance \stepcounter{utterance}  
    & & \multicolumn{4}{p{0.6\linewidth}}{
        \cellcolor[rgb]{0.9,0.9,0.9}{
            \makecell[{{p{\linewidth}}}]{
                \texttt{\tiny{[P1$\langle$GM]}}
                \texttt{Vārds jards tika uzminēts, bet tas ir nevainojams vārds. Jūsu komandas biedra gājiens beidzās tur.} \\
\texttt{Tagad sniedziet vēl vienu norādi, kas attiecas uz dažiem no jūsu atlikušajiem komandas vārdiem, un saistīto mērķa vārdu sarakstu. Neaizmirstiet sākt savu norādi ar 'PADOMS: ', pēc tam ievietot jaunu rindu un sākt ar komatiem atdalīto mērķa vārdu sarakstu ar 'MĒRĶI: '. Piezīme: salīdzinājumā ar iepriekšējiem pieprasījumiem daži vārdi no sarakstiem ir izņemti.} \\
\\ 
\texttt{Jūsu atlikušie komandas vārdi ir: vītītis, ātruma pārsniegšana, lietus, kardināls.} \\
\texttt{Atlikušie pretinieka vārdi ir: kliņģerīte, grāmatvedis, pistole, tops, ģerānija, mehāniķis.} \\
\texttt{Atlikušie traucētājvārdi ir: pēda, krekls, kaklasaite, sporta krekls, kilometrs, kostīms.} \\
\texttt{Atlikušie slepkavības vārdi ir: lielgabals.} \\
            }
        }
    }
    & \\ \\

    \theutterance \stepcounter{utterance}  
    & \multicolumn{4}{p{0.6\linewidth}}{
        \cellcolor[rgb]{0.9,0.9,0.9}{
            \makecell[{{p{\linewidth}}}]{
                \texttt{\tiny{[P1$\rangle$GM]}}
                \texttt{PADOMS: spirāle} \\
\texttt{MĒRĶI: vītītis, kardināls} \\
            }
        }
    }
    & & \\ \\

    \theutterance \stepcounter{utterance}  
    & & & \multicolumn{2}{p{0.3\linewidth}}{
        \cellcolor[rgb]{0.9,0.9,0.9}{
            \makecell[{{p{\linewidth}}}]{
                \texttt{\tiny{[GM$|$GM]}}
                \texttt{spirāle} \\
            }
        }
    }
    & & \\ \\

    \theutterance \stepcounter{utterance}  
    & & & \multicolumn{2}{p{0.3\linewidth}}{
        \cellcolor[rgb]{0.9,0.9,0.9}{
            \makecell[{{p{\linewidth}}}]{
                \texttt{\tiny{[GM$|$GM]}}
                ['vītītis', 'kardināls']
            }
        }
    }
    & & \\ \\

    \theutterance \stepcounter{utterance}  
    & & & \multicolumn{2}{p{0.3\linewidth}}{
        \cellcolor[rgb]{0.9,0.9,0.9}{
            \makecell[{{p{\linewidth}}}]{
                \texttt{\tiny{[GM$|$GM]}}
                {'word': 'vītītis', 'assignment': 'team'}
            }
        }
    }
    & & \\ \\

    \theutterance \stepcounter{utterance}  
    & & & \multicolumn{2}{p{0.3\linewidth}}{
        \cellcolor[rgb]{0.9,0.9,0.9}{
            \makecell[{{p{\linewidth}}}]{
                \texttt{\tiny{[GM$|$GM]}}
                {'word': 'kardināls', 'assignment': 'team'}
            }
        }
    }
    & & \\ \\

    \theutterance \stepcounter{utterance}  
    & & \multicolumn{4}{p{0.6\linewidth}}{
        \cellcolor[rgb]{0.9,0.9,0.9}{
            \makecell[{{p{\linewidth}}}]{
                \texttt{\tiny{[GM$\rangle$P2]}}
                \texttt{Vārds jards tika uzminēts, bet tas ir nevainojams vārds. Jūsu gājiens beidzās tur.} \\
\texttt{Tagad sniedziet vēl vienu ar komatiem atdalītu sarakstu, kurā ir vismaz 1 un ne vairāk kā 2 vārdi no turpmāk minētā vārdu saraksta, kas vislabāk sasaucas vai ir visciešāk saistīti ar vārdu 'spirāle'. Neaizmirstiet atbildi sākt ar MINĒJUMS: . Piezīme: salīdzinājumā ar iepriekšējiem pieprasījumiem daži vārdi no saraksta ir izņemti.} \\
\\ 
\texttt{lietus, vītītis, mehāniķis, kardināls, kliņģerīte, pistole, ģerānija, sporta krekls, kilometrs, kostīms, lielgabals, kaklasaite, grāmatvedis, ātruma pārsniegšana, krekls, pēda} \\
            }
        }
    }
    & \\ \\

    \theutterance \stepcounter{utterance}  
    & & & \multicolumn{4}{p{0.6\linewidth}}{
        \cellcolor[rgb]{0.9,0.9,0.9}{
            \makecell[{{p{\linewidth}}}]{
                \texttt{\tiny{[GM$\langle$P2]}}
                \texttt{MINĒJUMS: vītītis} \\
            }
        }
    }
     \\ \\

    \theutterance \stepcounter{utterance}  
    & & & \multicolumn{2}{p{0.3\linewidth}}{
        \cellcolor[rgb]{0.9,0.9,0.9}{
            \makecell[{{p{\linewidth}}}]{
                \texttt{\tiny{[GM$|$GM]}}
                ['vītītis']
            }
        }
    }
    & & \\ \\

    \theutterance \stepcounter{utterance}  
    & & & \multicolumn{2}{p{0.3\linewidth}}{
        \cellcolor[rgb]{0.9,0.9,0.9}{
            \makecell[{{p{\linewidth}}}]{
                \texttt{\tiny{[GM$|$GM]}}
                {'word': 'vītītis', 'assignment': 'team'}
            }
        }
    }
    & & \\ \\

    \theutterance \stepcounter{utterance}  
    & & & \multicolumn{2}{p{0.3\linewidth}}{
        \cellcolor[rgb]{0.9,0.9,0.9}{
            \makecell[{{p{\linewidth}}}]{
                \texttt{\tiny{[GM$|$GM]}}
                {'word': 'vītītis', 'assignment': 'team'}
            }
        }
    }
    & & \\ \\

    \theutterance \stepcounter{utterance}  
    & & & \multicolumn{2}{p{0.3\linewidth}}{
        \cellcolor[rgb]{0.9,0.9,0.9}{
            \makecell[{{p{\linewidth}}}]{
                \texttt{\tiny{[GM$|$GM]}}
                {'word': 'mehāniķis', 'assignment': 'opponent'}
            }
        }
    }
    & & \\ \\

    \theutterance \stepcounter{utterance}  
    & & & \multicolumn{2}{p{0.3\linewidth}}{
        \cellcolor[rgb]{0.9,0.9,0.9}{
            \makecell[{{p{\linewidth}}}]{
                \texttt{\tiny{[GM$|$GM]}}
                {'hidden': {'team': ['ātruma pārsniegšana', 'lietus', 'kardināls'], 'innocent': ['pēda', 'krekls', 'kaklasaite', 'sporta krekls', 'kilometrs', 'kostīms'], 'opponent': ['kliņģerīte', 'grāmatvedis', 'pistole', 'ģerānija'], 'assassin': ['lielgabals']}, 'revealed': {'team': {'team': ['taifūns', 'tornado', 'viesuļvētra', 'laupīšana', 'krāpšana', 'vītītis'], 'innocent': ['jards'], 'opponent': [], 'assassin': []}, 'opponent': {'team': [], 'innocent': [], 'opponent': ['jurists', 'roze', 'tops', 'mehāniķis'], 'assassin': []}}}
            }
        }
    }
    & & \\ \\

    \theutterance \stepcounter{utterance}  
    & & \multicolumn{4}{p{0.6\linewidth}}{
        \cellcolor[rgb]{0.9,0.9,0.9}{
            \makecell[{{p{\linewidth}}}]{
                \texttt{\tiny{[P1$\langle$GM]}}
                \texttt{Vārds vītītis tika uzminēts pareizi. Jūsu komandas biedra gājiens beidzās tur.} \\
\texttt{Tagad sniedziet vēl vienu norādi, kas attiecas uz dažiem no jūsu atlikušajiem komandas vārdiem, un saistīto mērķa vārdu sarakstu. Neaizmirstiet sākt savu norādi ar 'PADOMS: ', pēc tam ievietot jaunu rindu un sākt ar komatiem atdalīto mērķa vārdu sarakstu ar 'MĒRĶI: '. Piezīme: salīdzinājumā ar iepriekšējiem pieprasījumiem daži vārdi no sarakstiem ir izņemti.} \\
\\ 
\texttt{Jūsu atlikušie komandas vārdi ir: ātruma pārsniegšana, lietus, kardināls.} \\
\texttt{Atlikušie pretinieka vārdi ir: kliņģerīte, grāmatvedis, pistole, ģerānija, mehāniķis.} \\
\texttt{Atlikušie traucētājvārdi ir: pēda, krekls, kaklasaite, sporta krekls, kilometrs, kostīms.} \\
\texttt{Atlikušie slepkavības vārdi ir: lielgabals.} \\
            }
        }
    }
    & \\ \\

    \theutterance \stepcounter{utterance}  
    & \multicolumn{4}{p{0.6\linewidth}}{
        \cellcolor[rgb]{0.9,0.9,0.9}{
            \makecell[{{p{\linewidth}}}]{
                \texttt{\tiny{[P1$\rangle$GM]}}
                \texttt{PADOMS: baznīca} \\
\texttt{MĒRĶI: kardināls} \\
            }
        }
    }
    & & \\ \\

    \theutterance \stepcounter{utterance}  
    & & & \multicolumn{2}{p{0.3\linewidth}}{
        \cellcolor[rgb]{0.9,0.9,0.9}{
            \makecell[{{p{\linewidth}}}]{
                \texttt{\tiny{[GM$|$GM]}}
                \texttt{baznīca} \\
            }
        }
    }
    & & \\ \\

    \theutterance \stepcounter{utterance}  
    & & & \multicolumn{2}{p{0.3\linewidth}}{
        \cellcolor[rgb]{0.9,0.9,0.9}{
            \makecell[{{p{\linewidth}}}]{
                \texttt{\tiny{[GM$|$GM]}}
                ['kardināls']
            }
        }
    }
    & & \\ \\

    \theutterance \stepcounter{utterance}  
    & & & \multicolumn{2}{p{0.3\linewidth}}{
        \cellcolor[rgb]{0.9,0.9,0.9}{
            \makecell[{{p{\linewidth}}}]{
                \texttt{\tiny{[GM$|$GM]}}
                {'word': 'kardināls', 'assignment': 'team'}
            }
        }
    }
    & & \\ \\

    \theutterance \stepcounter{utterance}  
    & & \multicolumn{4}{p{0.6\linewidth}}{
        \cellcolor[rgb]{0.9,0.9,0.9}{
            \makecell[{{p{\linewidth}}}]{
                \texttt{\tiny{[GM$\rangle$P2]}}
                \texttt{Vārds vītītis tika uzminēts pareizi. Jūsu gājiens beidzās tur.} \\
\texttt{Tagad sniedziet vēl vienu ar komatiem atdalītu sarakstu, kurā ir vismaz 1 un ne vairāk kā 1 vārdi no turpmāk minētā vārdu saraksta, kas vislabāk sasaucas vai ir visciešāk saistīti ar vārdu 'baznīca'. Neaizmirstiet atbildi sākt ar MINĒJUMS: . Piezīme: salīdzinājumā ar iepriekšējiem pieprasījumiem daži vārdi no saraksta ir izņemti.} \\
\\ 
\texttt{lietus, kardināls, kliņģerīte, pistole, ģerānija, sporta krekls, kilometrs, kostīms, lielgabals, kaklasaite, grāmatvedis, ātruma pārsniegšana, krekls, pēda} \\
            }
        }
    }
    & \\ \\

    \theutterance \stepcounter{utterance}  
    & & & \multicolumn{4}{p{0.6\linewidth}}{
        \cellcolor[rgb]{0.9,0.9,0.9}{
            \makecell[{{p{\linewidth}}}]{
                \texttt{\tiny{[GM$\langle$P2]}}
                \texttt{MINĒJUMS: kardināls} \\
            }
        }
    }
     \\ \\

    \theutterance \stepcounter{utterance}  
    & & & \multicolumn{2}{p{0.3\linewidth}}{
        \cellcolor[rgb]{0.9,0.9,0.9}{
            \makecell[{{p{\linewidth}}}]{
                \texttt{\tiny{[GM$|$GM]}}
                ['kardināls']
            }
        }
    }
    & & \\ \\

    \theutterance \stepcounter{utterance}  
    & & & \multicolumn{2}{p{0.3\linewidth}}{
        \cellcolor[rgb]{0.9,0.9,0.9}{
            \makecell[{{p{\linewidth}}}]{
                \texttt{\tiny{[GM$|$GM]}}
                {'word': 'kardināls', 'assignment': 'team'}
            }
        }
    }
    & & \\ \\

    \theutterance \stepcounter{utterance}  
    & & & \multicolumn{2}{p{0.3\linewidth}}{
        \cellcolor[rgb]{0.9,0.9,0.9}{
            \makecell[{{p{\linewidth}}}]{
                \texttt{\tiny{[GM$|$GM]}}
                {'word': 'kardināls', 'assignment': 'team'}
            }
        }
    }
    & & \\ \\

    \theutterance \stepcounter{utterance}  
    & & & \multicolumn{2}{p{0.3\linewidth}}{
        \cellcolor[rgb]{0.9,0.9,0.9}{
            \makecell[{{p{\linewidth}}}]{
                \texttt{\tiny{[GM$|$GM]}}
                {'word': 'pistole', 'assignment': 'opponent'}
            }
        }
    }
    & & \\ \\

    \theutterance \stepcounter{utterance}  
    & & & \multicolumn{2}{p{0.3\linewidth}}{
        \cellcolor[rgb]{0.9,0.9,0.9}{
            \makecell[{{p{\linewidth}}}]{
                \texttt{\tiny{[GM$|$GM]}}
                {'hidden': {'team': ['ātruma pārsniegšana', 'lietus'], 'innocent': ['pēda', 'krekls', 'kaklasaite', 'sporta krekls', 'kilometrs', 'kostīms'], 'opponent': ['kliņģerīte', 'grāmatvedis', 'ģerānija'], 'assassin': ['lielgabals']}, 'revealed': {'team': {'team': ['taifūns', 'tornado', 'viesuļvētra', 'laupīšana', 'krāpšana', 'vītītis', 'kardināls'], 'innocent': ['jards'], 'opponent': [], 'assassin': []}, 'opponent': {'team': [], 'innocent': [], 'opponent': ['jurists', 'roze', 'tops', 'mehāniķis', 'pistole'], 'assassin': []}}}
            }
        }
    }
    & & \\ \\

    \theutterance \stepcounter{utterance}  
    & & \multicolumn{4}{p{0.6\linewidth}}{
        \cellcolor[rgb]{0.9,0.9,0.9}{
            \makecell[{{p{\linewidth}}}]{
                \texttt{\tiny{[P1$\langle$GM]}}
                \texttt{Vārds kardināls tika uzminēts pareizi. Jūsu komandas biedra gājiens beidzās tur.} \\
\texttt{Tagad sniedziet vēl vienu norādi, kas attiecas uz dažiem no jūsu atlikušajiem komandas vārdiem, un saistīto mērķa vārdu sarakstu. Neaizmirstiet sākt savu norādi ar 'PADOMS: ', pēc tam ievietot jaunu rindu un sākt ar komatiem atdalīto mērķa vārdu sarakstu ar 'MĒRĶI: '. Piezīme: salīdzinājumā ar iepriekšējiem pieprasījumiem daži vārdi no sarakstiem ir izņemti.} \\
\\ 
\texttt{Jūsu atlikušie komandas vārdi ir: ātruma pārsniegšana, lietus.} \\
\texttt{Atlikušie pretinieka vārdi ir: kliņģerīte, grāmatvedis, pistole, ģerānija.} \\
\texttt{Atlikušie traucētājvārdi ir: pēda, krekls, kaklasaite, sporta krekls, kilometrs, kostīms.} \\
\texttt{Atlikušie slepkavības vārdi ir: lielgabals.} \\
            }
        }
    }
    & \\ \\

    \theutterance \stepcounter{utterance}  
    & \multicolumn{4}{p{0.6\linewidth}}{
        \cellcolor[rgb]{0.9,0.9,0.9}{
            \makecell[{{p{\linewidth}}}]{
                \texttt{\tiny{[P1$\rangle$GM]}}
                \texttt{PADOMS: ūdens} \\
\texttt{MĒRĶI: lietus} \\
            }
        }
    }
    & & \\ \\

    \theutterance \stepcounter{utterance}  
    & & & \multicolumn{2}{p{0.3\linewidth}}{
        \cellcolor[rgb]{0.9,0.9,0.9}{
            \makecell[{{p{\linewidth}}}]{
                \texttt{\tiny{[GM$|$GM]}}
                \texttt{ūdens} \\
            }
        }
    }
    & & \\ \\

    \theutterance \stepcounter{utterance}  
    & & & \multicolumn{2}{p{0.3\linewidth}}{
        \cellcolor[rgb]{0.9,0.9,0.9}{
            \makecell[{{p{\linewidth}}}]{
                \texttt{\tiny{[GM$|$GM]}}
                ['lietus']
            }
        }
    }
    & & \\ \\

    \theutterance \stepcounter{utterance}  
    & & & \multicolumn{2}{p{0.3\linewidth}}{
        \cellcolor[rgb]{0.9,0.9,0.9}{
            \makecell[{{p{\linewidth}}}]{
                \texttt{\tiny{[GM$|$GM]}}
                {'word': 'lietus', 'assignment': 'team'}
            }
        }
    }
    & & \\ \\

    \theutterance \stepcounter{utterance}  
    & & \multicolumn{4}{p{0.6\linewidth}}{
        \cellcolor[rgb]{0.9,0.9,0.9}{
            \makecell[{{p{\linewidth}}}]{
                \texttt{\tiny{[GM$\rangle$P2]}}
                \texttt{Vārds kardināls tika uzminēts pareizi. Jūsu gājiens beidzās tur.} \\
\texttt{Tagad sniedziet vēl vienu ar komatiem atdalītu sarakstu, kurā ir vismaz 1 un ne vairāk kā 1 vārdi no turpmāk minētā vārdu saraksta, kas vislabāk sasaucas vai ir visciešāk saistīti ar vārdu 'ūdens'. Neaizmirstiet atbildi sākt ar MINĒJUMS: . Piezīme: salīdzinājumā ar iepriekšējiem pieprasījumiem daži vārdi no saraksta ir izņemti.} \\
\\ 
\texttt{lietus, kliņģerīte, ģerānija, sporta krekls, kilometrs, kostīms, lielgabals, kaklasaite, grāmatvedis, ātruma pārsniegšana, krekls, pēda} \\
            }
        }
    }
    & \\ \\

    \theutterance \stepcounter{utterance}  
    & & & \multicolumn{4}{p{0.6\linewidth}}{
        \cellcolor[rgb]{0.9,0.9,0.9}{
            \makecell[{{p{\linewidth}}}]{
                \texttt{\tiny{[GM$\langle$P2]}}
                \texttt{MINĒJUMS: lietus} \\
            }
        }
    }
     \\ \\

    \theutterance \stepcounter{utterance}  
    & & & \multicolumn{2}{p{0.3\linewidth}}{
        \cellcolor[rgb]{0.9,0.9,0.9}{
            \makecell[{{p{\linewidth}}}]{
                \texttt{\tiny{[GM$|$GM]}}
                ['lietus']
            }
        }
    }
    & & \\ \\

    \theutterance \stepcounter{utterance}  
    & & & \multicolumn{2}{p{0.3\linewidth}}{
        \cellcolor[rgb]{0.9,0.9,0.9}{
            \makecell[{{p{\linewidth}}}]{
                \texttt{\tiny{[GM$|$GM]}}
                {'word': 'lietus', 'assignment': 'team'}
            }
        }
    }
    & & \\ \\

    \theutterance \stepcounter{utterance}  
    & & & \multicolumn{2}{p{0.3\linewidth}}{
        \cellcolor[rgb]{0.9,0.9,0.9}{
            \makecell[{{p{\linewidth}}}]{
                \texttt{\tiny{[GM$|$GM]}}
                {'word': 'lietus', 'assignment': 'team'}
            }
        }
    }
    & & \\ \\

    \theutterance \stepcounter{utterance}  
    & & & \multicolumn{2}{p{0.3\linewidth}}{
        \cellcolor[rgb]{0.9,0.9,0.9}{
            \makecell[{{p{\linewidth}}}]{
                \texttt{\tiny{[GM$|$GM]}}
                {'word': 'kliņģerīte', 'assignment': 'opponent'}
            }
        }
    }
    & & \\ \\

    \theutterance \stepcounter{utterance}  
    & & & \multicolumn{2}{p{0.3\linewidth}}{
        \cellcolor[rgb]{0.9,0.9,0.9}{
            \makecell[{{p{\linewidth}}}]{
                \texttt{\tiny{[GM$|$GM]}}
                {'hidden': {'team': ['ātruma pārsniegšana'], 'innocent': ['pēda', 'krekls', 'kaklasaite', 'sporta krekls', 'kilometrs', 'kostīms'], 'opponent': ['grāmatvedis', 'ģerānija'], 'assassin': ['lielgabals']}, 'revealed': {'team': {'team': ['taifūns', 'tornado', 'viesuļvētra', 'laupīšana', 'krāpšana', 'vītītis', 'kardināls', 'lietus'], 'innocent': ['jards'], 'opponent': [], 'assassin': []}, 'opponent': {'team': [], 'innocent': [], 'opponent': ['jurists', 'roze', 'tops', 'mehāniķis', 'pistole', 'kliņģerīte'], 'assassin': []}}}
            }
        }
    }
    & & \\ \\

    \theutterance \stepcounter{utterance}  
    & & \multicolumn{4}{p{0.6\linewidth}}{
        \cellcolor[rgb]{0.9,0.9,0.9}{
            \makecell[{{p{\linewidth}}}]{
                \texttt{\tiny{[P1$\langle$GM]}}
                \texttt{Vārds lietus tika uzminēts pareizi. Jūsu komandas biedra gājiens beidzās tur.} \\
\texttt{Tagad sniedziet vēl vienu norādi, kas attiecas uz dažiem no jūsu atlikušajiem komandas vārdiem, un saistīto mērķa vārdu sarakstu. Neaizmirstiet sākt savu norādi ar 'PADOMS: ', pēc tam ievietot jaunu rindu un sākt ar komatiem atdalīto mērķa vārdu sarakstu ar 'MĒRĶI: '. Piezīme: salīdzinājumā ar iepriekšējiem pieprasījumiem daži vārdi no sarakstiem ir izņemti.} \\
\\ 
\texttt{Jūsu atlikušie komandas vārdi ir: ātruma pārsniegšana.} \\
\texttt{Atlikušie pretinieka vārdi ir: kliņģerīte, grāmatvedis, ģerānija.} \\
\texttt{Atlikušie traucētājvārdi ir: pēda, krekls, kaklasaite, sporta krekls, kilometrs, kostīms.} \\
\texttt{Atlikušie slepkavības vārdi ir: lielgabals.} \\
            }
        }
    }
    & \\ \\

    \theutterance \stepcounter{utterance}  
    & \multicolumn{4}{p{0.6\linewidth}}{
        \cellcolor[rgb]{0.9,0.9,0.9}{
            \makecell[{{p{\linewidth}}}]{
                \texttt{\tiny{[P1$\rangle$GM]}}
                \texttt{PADOMS: sods} \\
\texttt{MĒRĶI: ātruma pārsniegšana} \\
            }
        }
    }
    & & \\ \\

    \theutterance \stepcounter{utterance}  
    & & & \multicolumn{2}{p{0.3\linewidth}}{
        \cellcolor[rgb]{0.9,0.9,0.9}{
            \makecell[{{p{\linewidth}}}]{
                \texttt{\tiny{[GM$|$GM]}}
                \texttt{sods} \\
            }
        }
    }
    & & \\ \\

    \theutterance \stepcounter{utterance}  
    & & & \multicolumn{2}{p{0.3\linewidth}}{
        \cellcolor[rgb]{0.9,0.9,0.9}{
            \makecell[{{p{\linewidth}}}]{
                \texttt{\tiny{[GM$|$GM]}}
                ['ātruma pārsniegšana']
            }
        }
    }
    & & \\ \\

    \theutterance \stepcounter{utterance}  
    & & & \multicolumn{2}{p{0.3\linewidth}}{
        \cellcolor[rgb]{0.9,0.9,0.9}{
            \makecell[{{p{\linewidth}}}]{
                \texttt{\tiny{[GM$|$GM]}}
                {'word': 'ātruma pārsniegšana', 'assignment': 'team'}
            }
        }
    }
    & & \\ \\

    \theutterance \stepcounter{utterance}  
    & & \multicolumn{4}{p{0.6\linewidth}}{
        \cellcolor[rgb]{0.9,0.9,0.9}{
            \makecell[{{p{\linewidth}}}]{
                \texttt{\tiny{[GM$\rangle$P2]}}
                \texttt{Vārds lietus tika uzminēts pareizi. Jūsu gājiens beidzās tur.} \\
\texttt{Tagad sniedziet vēl vienu ar komatiem atdalītu sarakstu, kurā ir vismaz 1 un ne vairāk kā 1 vārdi no turpmāk minētā vārdu saraksta, kas vislabāk sasaucas vai ir visciešāk saistīti ar vārdu 'sods'. Neaizmirstiet atbildi sākt ar MINĒJUMS: . Piezīme: salīdzinājumā ar iepriekšējiem pieprasījumiem daži vārdi no saraksta ir izņemti.} \\
\\ 
\texttt{ģerānija, sporta krekls, kilometrs, kostīms, lielgabals, kaklasaite, grāmatvedis, ātruma pārsniegšana, krekls, pēda} \\
            }
        }
    }
    & \\ \\

    \theutterance \stepcounter{utterance}  
    & & & \multicolumn{4}{p{0.6\linewidth}}{
        \cellcolor[rgb]{0.9,0.9,0.9}{
            \makecell[{{p{\linewidth}}}]{
                \texttt{\tiny{[GM$\langle$P2]}}
                \texttt{MINĒJUMS: ātruma pārsniegšana} \\
            }
        }
    }
     \\ \\

    \theutterance \stepcounter{utterance}  
    & & & \multicolumn{2}{p{0.3\linewidth}}{
        \cellcolor[rgb]{0.9,0.9,0.9}{
            \makecell[{{p{\linewidth}}}]{
                \texttt{\tiny{[GM$|$GM]}}
                ['ātruma pārsniegšana']
            }
        }
    }
    & & \\ \\

    \theutterance \stepcounter{utterance}  
    & & & \multicolumn{2}{p{0.3\linewidth}}{
        \cellcolor[rgb]{0.9,0.9,0.9}{
            \makecell[{{p{\linewidth}}}]{
                \texttt{\tiny{[GM$|$GM]}}
                {'word': 'ātruma pārsniegšana', 'assignment': 'team'}
            }
        }
    }
    & & \\ \\

    \theutterance \stepcounter{utterance}  
    & & & \multicolumn{2}{p{0.3\linewidth}}{
        \cellcolor[rgb]{0.9,0.9,0.9}{
            \makecell[{{p{\linewidth}}}]{
                \texttt{\tiny{[GM$|$GM]}}
                {'word': 'ātruma pārsniegšana', 'assignment': 'team'}
            }
        }
    }
    & & \\ \\

    \theutterance \stepcounter{utterance}  
    & & & \multicolumn{2}{p{0.3\linewidth}}{
        \cellcolor[rgb]{0.9,0.9,0.9}{
            \makecell[{{p{\linewidth}}}]{
                \texttt{\tiny{[GM$|$GM]}}
                \texttt{team has won} \\
            }
        }
    }
    & & \\ \\

\end{supertabular}
}

\end{document}
